% !PRINTED 529
\weberpage{529}
%<<CONT
und erhalten also, da nach (10)
\weq{15}{\varphi_+(0) = \varphi_-(0) = \varphi(0) = \int_0^\infty e^{\pi iyv(\omega_1+\omega_2)}\left(\frac{v}{2}\right)^{2s-2}dv}
ist, aus (11), (12) und (14):
\weq{16}{\Sigma_{-\infty,\infty}^{x}W = \varphi(0) + \Sigma_{1,\infty}^{\nu}[\varphi_+(\nu)+\varphi_-(\nu)].}

Führen wir dies in (8) ein, so zerfällt $M_s$ in zwei Teile:
\weq{17}{M_s = P_s + Q_s,}
wenn
\weq{18}{P_s = \frac{(2\pi)^{2s}}{A^s\Gamma(s)\Gamma(s)}\Sigma_{1,\infty}^{y}\varphi(0),}
\weq{19}{Q_s = \frac{(2\pi)^{2s}}{A^s\Gamma(s)\Gamma(s)}\Sigma_{1,\infty}^{y}\Sigma_{1,\infty}^{\nu}[\varphi_+(\nu)+\varphi_-(\nu)]}
gesetzt wird.

Nach (15) ist, wenn wir für $v$ eine neue Integrationsvariable
$t$ durch die Gleichung
\[ \pi iv(\omega_1+\omega_2) = -\frac{2\pi v\sqrt{m}}{A} = -t \]
einführen, und die Summation nach $y$ ausführen:
\[ \Sigma_{1,\infty}^{y}\varphi(0) = \frac{2A^{2s-1}}{(4\pi\sqrt{m})^{2s-1}}\int_0^\infty\frac{t^{2s-2}\,dt}{e^t-1}, \]
also:
\weq{20}{P_s = \frac{(2\pi)^{2s}A^{s-1}}{(4\pi\sqrt{m})^{2s-1}\Gamma(s)^2}\,2\int_0^\infty\frac{t^{2s-2}\,dt}{e^t-1}.}

Hieraus läßt sich der Grenzwert $P$ leicht bestimmen. Es
ist nämlich das Integral
\weq{21}{\int_0^\infty t^{2s-2}e^{-t}\left(\frac{1}{1-e^{-t}}-\frac{1}{t}\right)dt,}
da der in der Klammer stehende Ausdruck für $t=0$ und $t=\infty$
einen endlichen Wert behält, bis $s=1$ stetig, und hat daher
den Grenzwert:
\[ \int_0^\infty e^{-t}\left(\frac{1}{1-e^{-t}}-\frac{1}{t}\right)dt = -\Gamma'(1) = 0{,}5772\ldots \]
